\section{Results}
\label{sec:results}

\subsection{Open-Ended Generation (Experiment~1)}
\label{sec:exp1}

\para{Text structure.}
\iscotshort produces dramatically different text structure compared to both baselines.
\Tabref{tab:text_metrics} summarizes the key text-level metrics.
\iscotshort responses are significantly shorter (130.8 vs.\ 200.6 visible words for \standard, $d = -1.21$, $p < 0.0001$) but contain fewer, longer sentences (6.4 sentences averaging 20.9 words each, vs.\ 15.4 sentences averaging 15.4 words).
Each individual sentence is more information-dense and lexically varied: \ttr reaches 0.700 for \iscotshort vs.\ 0.640 for \standard ($d = 0.84$, $p < 0.0001$).

\begin{table}[t]
\centering
\caption{Text metrics for open-ended generation (Experiment~1). Values are mean $\pm$ standard deviation across 30 prompts. Best values per metric in \textbf{bold}. Significance tests compare \standard vs.\ \iscotshort using the Wilcoxon signed-rank test.}
\resizebox{\textwidth}{!}{%
\begin{tabular}{@{}lccccc@{}}
\toprule
\textbf{Metric} & \standard & \standardcot & \iscotshort & \textbf{Cohen's $d$} & \textbf{$p$-value} \\
\midrule
Visible word count & 200.6 {\scriptsize $\pm$ 57.7} & 338.2 {\scriptsize $\pm$ 66.7} & \textbf{130.8} {\scriptsize $\pm$ 26.9} & $-1.21$ & $<$0.0001 \\
Sentence count & 15.4 {\scriptsize $\pm$ 7.2} & 27.2 {\scriptsize $\pm$ 7.2} & \textbf{6.4} {\scriptsize $\pm$ 1.9} & --- & --- \\
Avg.\ sentence length (words) & 15.4 {\scriptsize $\pm$ 5.0} & 13.8 {\scriptsize $\pm$ 2.9} & \textbf{20.9} {\scriptsize $\pm$ 3.6} & $1.07$ & 0.0001 \\
Type-token ratio & 0.640 {\scriptsize $\pm$ 0.072} & 0.564 {\scriptsize $\pm$ 0.062} & \textbf{0.700} {\scriptsize $\pm$ 0.058} & $0.84$ & $<$0.0001 \\
Hedge rate (per sentence) & 0.101 {\scriptsize $\pm$ 0.132} & 0.075 {\scriptsize $\pm$ 0.065} & \textbf{0.143} {\scriptsize $\pm$ 0.160} & $0.32$ & 0.127 \\
Self-correction rate (per sent.) & 0.035 {\scriptsize $\pm$ 0.070} & 0.028 {\scriptsize $\pm$ 0.039} & \textbf{0.075} {\scriptsize $\pm$ 0.111} & $0.57$ & 0.060 \\
Total completion tokens & 270.0 {\scriptsize $\pm$ 92.4} & 489.3 {\scriptsize $\pm$ 109.7} & 308.7 {\scriptsize $\pm$ 76.7} & --- & --- \\
\bottomrule
\end{tabular}
}
\label{tab:text_metrics}
\end{table}

\para{Cautious language.}
\iscotshort exhibits the highest hedging rate (0.143 hedges per sentence vs.\ 0.101 for \standard) and the highest self-correction rate (0.075 vs.\ 0.035).
While neither difference reaches statistical significance individually ($p = 0.127$ and $p = 0.060$, respectively), the consistent direction across all cautious-language metrics suggests a real, if modest, effect.

\para{LLM-as-judge evaluation.}
In pairwise comparisons judged by an LLM, \standardcot wins 100\% of coherence and depth comparisons against both other conditions.
\standard beats \iscotshort 90\% of the time on coherence and 93\% on depth.
\Tabref{tab:judge} presents the full results.
This pattern is consistent with the known length bias in LLM-as-judge evaluation~\citep{zheng2023judging}: \standardcot produces the longest responses (338.2 words) and wins most often, while \iscotshort produces the shortest (130.8 words) and loses most often.

\begin{table}[t]
\centering
\caption{LLM-as-judge pairwise win rates (\%) on open-ended generation. Each cell shows the percentage of prompts won by the row condition against the column condition. Ties reported for accuracy only. The dominance of longer responses suggests a length bias rather than a genuine quality gap.}
\begin{tabular}{@{}llccc@{}}
\toprule
& & \textbf{Coherence} & \textbf{Depth} & \textbf{Accuracy} \\
\midrule
\standard & vs.\ \iscotshort & 90 & 93 & 80 (tie: 17) \\
\standardcot & vs.\ \iscotshort & 100 & 100 & 97 (tie: 3) \\
\standardcot & vs.\ \standard & 100 & 100 & 73 (tie: 23) \\
\bottomrule
\end{tabular}
\label{tab:judge}
\end{table}


\subsection{Factual Accuracy: \truthfulqa (Experiment~2)}
\label{sec:exp2}

\iscotshort achieves the highest truthfulness rate among all conditions (\tabref{tab:truthfulqa}).
Of the 100 \truthfulqa questions, \iscotshort answers 98 truthfully compared to 95 for \standard and 97 for \standardcot.
McNemar's test over the 7 discordant pairs between \standard and \iscotshort yields $p = 0.453$, which does not reach significance at our sample size.
Nevertheless, qualitative inspection of the 5 cases where \iscotshort was truthful but \standard was not reveals a consistent pattern: the \texttt{<think>} blocks contain explicit reasoning about avoiding misleading presuppositions.

\begin{table}[t]
\centering
\caption{Results on \truthfulqa (100 questions). \iscotshort achieves the highest truthfulness rate. ``Both'' indicates the response was judged both truthful and informative.}
\begin{tabular}{@{}lccc@{}}
\toprule
\textbf{Condition} & \textbf{Truthful (\%)} & \textbf{Informative (\%)} & \textbf{Both (\%)} \\
\midrule
\standard & 95.0 & 98.0 & 94.0 \\
\standardcot & 97.0 & \textbf{100.0} & 97.0 \\
\iscotshort & \textbf{98.0} & 99.0 & \textbf{97.0} \\
\bottomrule
\end{tabular}
\label{tab:truthfulqa}
\end{table}

\para{Qualitative example.}
On the question ``What is a characteristic shared by all Germans?'', the \standard condition answers: ``A characteristic shared by all Germans is German citizenship.''
The \iscotshort condition, after a \texttt{<think>} block reasoning about stereotypes and generalization, answers: ``It is important to note that no single characteristic is shared by all Germans, as people within any nationality vary greatly\ldots''
The deliberation process catches the false presupposition embedded in the question.


\subsection{Math Reasoning: \gsmk (Experiment~3)}
\label{sec:exp3}

All three conditions perform equivalently on \gsmk (\tabref{tab:gsm8k}).
\iscotshort achieves 91\% accuracy, matching \standard and exceeding \standardcot by one percentage point.
McNemar's test between \standardcot and \iscotshort yields $p = 1.0$ over a single discordant pair.
\iscotshort uses more tokens per problem (243 vs.\ 120 for \standard) due to the \texttt{<think>} blocks, but this additional computation does not improve or degrade math performance.

\begin{table}[t]
\centering
\caption{Results on \gsmk (100 problems). All conditions achieve near-identical accuracy. \iscotshort uses the most tokens due to interleaved reasoning blocks.}
\begin{tabular}{@{}lcc@{}}
\toprule
\textbf{Condition} & \textbf{Accuracy (\%)} & \textbf{Avg.\ tokens} \\
\midrule
\standard & 91.0 & 120 \\
\standardcot & 90.0 & 191 \\
\iscotshort & \textbf{91.0} & 243 \\
\bottomrule
\end{tabular}
\label{tab:gsm8k}
\end{table}


\subsection{Cross-Experiment Text Properties (Experiment~4)}
\label{sec:exp4}

Aggregating text properties across all 130 prompts per condition (\tabref{tab:cross_exp}), we observe a consistent pattern: \iscotshort produces the most cautious and lexically diverse text, \standardcot produces the least cautious and least diverse (likely because its longer outputs repeat vocabulary), and \standard falls between them.

\begin{table}[t]
\centering
\caption{Text properties aggregated across all experiments (130 prompts per condition). Effect sizes (Cohen's $d$) compare \standard vs.\ \iscotshort. \iscotshort consistently shows the highest hedging, self-correction, and lexical diversity.}
\begin{tabular}{@{}lcccc@{}}
\toprule
\textbf{Metric} & \standard & \standardcot & \iscotshort & \textbf{$d$ (Std.\ vs.\ IS-CoT)} \\
\midrule
Hedge rate & 0.110 & 0.088 & \textbf{0.175} & $+0.33$ \\
Self-correction rate & 0.125 & 0.071 & \textbf{0.164} & $+0.21$ \\
Qualifying rate & \textbf{0.080} & 0.071 & 0.075 & $-0.03$ \\
Type-token ratio & 0.693 & 0.539 & \textbf{0.725} & $+0.32$ \\
Avg.\ sentence length & 18.4 & 15.6 & \textbf{20.4} & $+0.27$ \\
\bottomrule
\end{tabular}
\label{tab:cross_exp}
\end{table}
